% ============================================================================
\chapter{Complete blueprint-to-module map}\label{ch:lean-files}
% ============================================================================
\module{(this chapter is the answer)}

This chapter records the complete assignment of all 378 target and context
items to 124 planned Lean modules in the existing \texttt{no\_compromise}
project.  The module root is \texttt{NoCompromise}; mathematical declarations
retain the \texttt{LiquidDrop} namespace.  All 302 mathematical results have
an owner, and all 59 remarks and formalisation notes are attached separately
as guidance.  No range or wildcard stands for an unassigned item.

The tree below lists each module with the items it owns.  The assignment
was checked mechanically for coverage and acyclicity using every declared
\verb|\uses| edge, shared conventions, and documented context prerequisites
(including the regression tests' references).  This is a plan for Lean
files, not an assertion that their statements or proofs have been elaborated.
Additional supporting lemmas may be needed during formalisation.

\section{Design rules used}

\begin{enumerate}[label=(R\arabic*),leftmargin=*]
\item\label{r:one-concept} One mathematical concept per file.  A file that
  needs a section heading to be navigable should be split.
\item\label{r:leaf-imports} Planned imports form a DAG.  Build in a
  topological order of the module imports; the thematic layers of
  Chapter~\ref{ch:dependencies} are only an overview.
\item\label{r:interface} Each layer has a single \texttt{...Basic} or
  \texttt{...Defs} file holding its definitions, so that downstream files
  import definitions without importing proofs.
\item\label{r:generic} The regularity theory is stated for
  \texttt{OmegaMinimal} (Definition~\ref{def:omega-minimal}) and never for the
  liquid-drop energy, per Remark~\ref{rem:generic-discipline}.  The energy
  appears only in \texttt{Energy/}, \code{Regularity/Penalization.lean}, and
  \texttt{Nonexistence/}, \texttt{Classification/}, \texttt{Binding/}.
\item\label{r:conventions} Each convention has one owner at the layer where
  its objects are available.  \texttt{Conventions.lean} holds the primitive
  ambient setting, null-set convention, test classes and constant discipline.
  The polar sign belongs to \texttt{BV/Basic.lean}; curvature belongs to
  \texttt{Surface/Geometry.lean}; the first-variation sign belongs with its
  theorems; the density representative belongs to the regularity layer.
  Mathematical claims inside a convention must still be proved.
\item\label{r:constants} Every quantitative statement takes its constants as
  explicit hypotheses or returns them existentially in the declared order of
  Convention~\ref{conv:constants}.  Constants are never section variables.
\item\label{r:no-heavy-defs} Files in the \texttt{Sard/}, \texttt{Flow/},
  \texttt{Topology/} and \texttt{Surface/} trees import nothing from
  \texttt{BV/} in the planned graph, so they can be built independently of the whole
  measure-theoretic half.  \texttt{Surface/} does import \texttt{Area/}
  (for \texttt{thm:area-formula-rect}, \texttt{lem:sphere-area} and
  \code{lem:smooth-surface-rectifiable}), which itself depends only on the
  base.
\end{enumerate}

\section{The tree}

% Lean module tree: each file with the blueprint items it owns.
{\small
\begin{verbatim}
NoCompromise.lean  -- root: imports Main and Binding.Main

NoCompromise/

  Area/
    AlmostLinear.lean
      -- const:almost-linear
      -- lem:almost-linear
    Cofactor.lean
      -- lem:cofactor
    Formula.lean
      -- thm:area-formula-2d
    GoodPieces.lean
      -- def:good-pieces
    GoodPiecesCover.lean
      -- lem:good-pieces-exhaust
    Graph.lean
      -- cor:graph-area
    Linear.lean
      -- def:jacobian
      -- lem:linear-distortion
    PieceFormula.lean
      -- lem:area-one-piece
    RankDeficient.lean
      -- lem:rank-deficient
    Rectifiable.lean
      -- def:rectifiable
      -- thm:area-formula-rect
      -- lem:smooth-surface-rectifiable
    Sphere.lean
      -- lem:sphere-area

  BV/
    AnnularGluingFull.lean
      -- prop:annular-gluing
    Basic.lean
      -- conv:outward
      -- lem:test-class
      -- thm:polar
      -- lem:perimeter-invariances
      -- lem:lsc
      -- lem:locality
    BoundaryTraces.lean
      -- thm:traces
    Coarea.lean
      -- thm:bv-coarea
    CoareaCritical.lean
      -- lem:coarea-critical-part
    CoareaL1Lebesgue.lean
      -- cor:coarea-L1
    CoareaRegular.lean
      -- lem:coarea-regular-part
    Compactness.lean
      -- lem:translation-estimate
      -- lem:grid-approx
      -- thm:bv-compactness
    Defs.lean
      -- def:density
      -- def:variation
      -- def:perimeter
    ExactCuts.lean
      -- prop:cut-identities
    Gluing.lean
      -- prop:gluing
    GoodRadii.lean
      -- def:good-radius
      -- lem:good-radius-ae
    GoodTruncation.lean
      -- lem:good-truncation
    JumpDisintegration.lean
      -- lem:bv-slices
    Rellich.lean
      -- lem:L2-interp
    Slicing.lean
      -- cor:slicing
    SmoothApprox.lean
      -- thm:smooth-approx
      -- cor:smooth-approx-energy
    StrictApprox.lean
      -- lem:bv-product-smooth
      -- lem:strict-approx
    WeightedCoarea.lean
      -- thm:weighted-coarea

  Ball/
    Defs.lean
      -- not:BV
    Perimeter.lean
      -- lem:ball-perimeter
    Potential.lean
      -- lem:ball-coulomb
      -- cor:ball-energy

  Binding/
    BallRatio.lean
      -- lem:optimal-ball-volume
    Main.lean
      -- thm:main-binding
    Ratio.lean
      -- def:relaxed-ratio
      -- lem:relaxed-attained
      -- lem:ratio-is-fixed
      -- lem:stabilization
    Splitting.lean
      -- lem:scaling-lower
      -- def:splitting-functions
      -- lem:ND-positive
      -- lem:two-piece
      -- lem:splitting-substitution
      -- lem:splitting-F
      -- thm:splitting
    Strict.lean
      -- prop:strict-binding

  CapacitaryK/
    Asymptotics.lean
      -- conv:level-orientation
      -- lem:K-normalization
      -- lem:K-level-asymptotics
      -- def:K-F
      -- lem:K-far-field
    Inequality.lean
      -- lem:K-gauss-bonnet-input
      -- prop:K-measure-inequality
      -- lem:K-h-monotone
      -- prop:K-two-ineq
      -- thm:capacitary-inequalities
    PositiveMeasure.lean
      -- lem:K-bochner
      -- prop:K-mu
      -- lem:K-level-identities
    Representatives.lean
      -- def:K-p
      -- lem:K-p-geometric
      -- def:K-Fhat
      -- prop:K-structure
      -- lem:K-pushforward-density
    SlabIdentities.lean
      -- lem:K-slab-H
      -- lem:K-slab-F
      -- lem:gradient-vanishes-on-zero-set

  Capacity/
    Flux.lean
      -- def:capacity
      -- lem:potential-decay
      -- lem:flux-identity
    Kelvin.lean
      -- lem:kelvin
      -- lem:level-asymptotics
    Levels.lean
      -- lem:level-connected
    Potential.lean
      -- thm:capacitary-potential
      -- not:w

  Classification/
    Endpoint.lean
      -- prop:endpoint
    Subcritical.lean
      -- not:z-delta
      -- lem:one-ball-upper
      -- prop:classification

  Compactness/
    Decomposition.lean
      -- thm:decomposition
      -- lem:radii-construction
      -- lem:perimeter-splitting
      -- lem:coulomb-splitting
    Nonvanishing.lean
      -- lem:cube-nonvanishing
      -- lem:coulomb-lsc
      -- thm:compactness

  Cones/
    NoSingular.lean
      -- prop:no-singular-points
    ThreeDim.lean
      -- lem:cone-tangent-cylinder
      -- lem:cone-descent
      -- lem:cone-smooth
      -- lem:cone-link-great-circle
      -- thm:cone-3d
    TwoDim.lean
      -- lem:cone-2d-link
      -- lem:cone-2d
  Conventions.lean
    -- conv:ambient
    -- conv:null
    -- conv:test-fields
    -- conv:constants

  DeGiorgi/
    Blowup.lean
      -- thm:halfspace-blowup
    ConeConcentration.lean
      -- lem:cone-concentration
    ExactDensity.lean
      -- thm:exact-density
    NormalFlux.lean
      -- lem:normal-flux
    Rectifiability.lean
      -- lem:rectifiability-criterion
    Reduced.lean
      -- not:degiorgi
      -- def:reduced-boundary
      -- lem:polar-differentiation
    ReducedDensity.lean
      -- lem:reduced-density
    ReducedPerimeter.lean
      -- lem:reduced-perimeter
    SmoothBoundary.lean
      -- cor:smooth-perimeter
    Structure.lean
      -- thm:degiorgi-structure
      -- thm:gauss-green

  Elliptic/
    Boundary.lean
      -- not:half-ball
      -- thm:boundary-C1a
      -- thm:boundary-C2a
      -- thm:boundary-nondiv
      -- thm:boundary-neumann
    Caccioppoli.lean
      -- lem:caccioppoli
    CampanatoComparison.lean
      -- lem:campanato-compare
    CampanatoGrowth.lean
      -- lem:campanato-energy-growth
    CampanatoHolder.lean
      -- thm:campanato
    CampanatoHolderDatum.lean
      -- lem:Gh
    CampanatoIteration.lean
      -- lem:campanato-iteration
    ClassicalAdmissible.lean
      -- thm:W11-gauss-green
    ClassicalKernelFlux.lean
      -- cor:W11-GG-uses
    FrozenDecay.lean
      -- lem:frozen-decay
    HarmonicDerivative.lean
      -- lem:harmonic-deriv
    HarmonicMeanValue.lean
      -- lem:mean-value
    HolderInterpolation.lean
      -- lem:holder-interp
    Hopf.lean
      -- prop:hopf
    HopfExterior.lean
      -- cor:hopf-exterior
    InteriorH2.lean
      -- lem:interior-H2
    KelvinRemovable.lean
      -- prop:kelvin-removable
    Newtonian.lean
      -- prop:newtonian
    NewtonianSchauder.lean
      -- thm:newtonian-schauder
    NondivSchauder.lean
      -- thm:nondiv-schauder
    Quasilinear.lean
      -- thm:quasilinear-schauder
    Schauder.lean
      -- const:nondiv-order
    SobolevChain.lean
      -- lem:sobolev-Ck
    StrongMaximumHarmonic.lean
      -- prop:strong-max
    TraceL1.lean
      -- conv:domain-classes
      -- prop:trace-L1
    WeakDirichlet.lean
      -- prop:weak-dirichlet
    WeakMaximum.lean
      -- lem:weak-max
    WeakSolutions.lean
      -- prop:weak-neumann

  Energy/
    Coulomb.lean
      -- not:C0
      -- lem:potential-bound
      -- lem:coulomb-bilinear
      -- lem:coulomb-lipschitz
      -- lem:coulomb-disjoint
      -- lem:cross-decay
    CoulombDefs.lean
      -- def:coulomb
    Defs.lean
      -- def:energy
      -- def:minimizer
    NullInvariance.lean
      -- lem:null-invariance
    PotentialHolder.lean
      -- lem:potential-C1beta
    Scaling.lean
      -- lem:coulomb-scaling

  Existence/
    Subcritical.lean
      -- prop:existence-subcritical

  Flow/
    FlowBox.lean
      -- thm:flow-box
      -- cor:flow-manifold
    ODE.lean
      -- thm:global-flow
      -- lem:flow-group
      -- lem:flow-gronwall
      -- lem:flow-variational
      -- thm:flow-Ck

  Green/
    Identity.lean
      -- not:v-omega
      -- lem:interior-flux
      -- lem:v-expansion
      -- lem:wronskian
      -- thm:green-identity

  Hull/
    Defs.lean
      -- def:hull
    Properties.lean
      -- lem:hull-properties

  Isoperimetric/
    ABP.lean
      -- lem:abp-neumann
      -- lem:abp-contact
      -- prop:iso-smooth
      -- lem:concavity-23
      -- cor:iso-smooth-components
      -- thm:sharp-isoperimetric
    Rigidity.lean
      -- thm:isoperimetric-rigidity
  Main.lean
    -- thm:main

  Measure/
    SignedRiesz.lean
      -- def:locally-bounded
      -- def:upper-variation
      -- lem:upper-variation-additive
      -- thm:signed-riesz
      -- lem:riesz-exhaustion
    WeakStar.lean
      -- lem:C0-separable
      -- def:exhaustion-cutoffs
      -- thm:weak-star-compactness

  Nonexistence/
    Sharp.lean
      -- not:L
      -- lem:two-ball-comparison
      -- prop:nonexistence-6
      -- prop:nonexistence-8
      -- cor:nonexistence
    Slicing.lean
      -- lem:slicing-inequality
      -- lem:slicing-integrals
      -- prop:V-le-8

  Regularity/
    ApproxHarmonic.lean
      -- lem:approx-harmonic
    CompressionJacobian.lean
      -- lem:compression-jacobian
    CompressionProfileTheorem.lean
      -- lem:compression-profile
    Cylinders.lean
      -- not:cylinders
    Deformation.lean
      -- thm:deformation
    Density.lean
      -- const:rd
    DensityAhlfors.lean
      -- lem:ahlfors
    DensityCutComparison.lean
      -- lem:cut-comparison
    DensityEstimates.lean
      -- lem:two-sided-density
    DensityRadial.lean
      -- not:mx-nx
      -- lem:mx-derivative
    EpsReg.lean
      -- lem:normals-cauchy
      -- lem:uniform-centers
      -- thm:eps-regularity
    Excess.lean
      -- def:excess
      -- lem:excess-ball-cyl
      -- lem:excess-center
    FirstVariation.lean
      -- lem:bounded-H
    FixedNormalExcess.lean
      -- lem:fixed-normal-excess-convergence
    FluxDefect.lean
      -- lem:flux-defect
    GeometricRecurrence.lean
      -- lem:geometric-recurrence
    GraphApprox.lean
      -- thm:graph-approx
    GraphBadBase.lean
      -- lem:bad-base-area
    GraphGoodBase.lean
      -- lem:good-base
    GraphProjectedExcess.lean
      -- def:projected-excess
    GraphSingleValued.lean
      -- lem:single-valued
    GraphTwoPointCaps.lean
      -- const:graph
      -- lem:two-point-lipschitz
    HarmonicAffine.lean
      -- lem:harmonic-affine
    HarmonicBlowup.lean
      -- lem:blowup-compactness
    HarmonicBlowupUniform.lean
      -- lem:harmonic-approx-uniform
    HeightBound.lean
      -- const:height
      -- lem:height-bound
    Monotonicity.lean
      -- thm:monotonicity
    OmegaMinimal.lean
      -- def:omega-minimal
    Penalization.lean
      -- lem:penalization
    PenalizationDilation.lean
      -- lem:dilation-quotients
    PenalizationQuasiminimal.lean
      -- lem:quasiminimal
    PerimeterConvergence.lean
      -- lem:perimeter-measure-convergence
    RadialDivergence.lean
      -- lem:radial-divergence
    RepresentativeBoundary.lean
      -- conv:representative
      -- lem:open-representative
    RepresentativeBounded.lean
      -- lem:bounded-representative
    ReversePoincare.lean
      -- lem:reverse-poincare
    SlabCap.lean
      -- def:slab-cap
    TangentCone.lean
      -- prop:tangent-is-cone
    TangentExistence.lean
      -- lem:tangent-cone-minimizing
    TiltImprovement.lean
      -- const:tilt
      -- lem:tilt-improvement
      -- cor:excess-decay
    ZeroExcess.lean
      -- lem:zero-excess

  Sard/
    Equidimensional.lean
      -- lem:sard-equidim
      -- cor:sard-charts
    FourToThree.lean
      -- thm:sard-4to3
    OneDimensional.lean
      -- lem:sard-1d
    Scalar.lean
      -- lem:sard-2d
    ThreeDimensional.lean
      -- thm:sard-3d

  Sobolev/
    BV.lean
      -- lem:GN-smooth
      -- thm:bv-sobolev
      -- lem:planar-sobolev-L2
    Extension.lean
      -- lem:bilip-chain
      -- lem:bv-extension
    H1Extension.lean
      -- lem:H1-extension-disk
    Maximal.lean
      -- def:maximal
      -- lem:maximal
    PlanarDomain.lean
      -- lem:planar-GN
    Poincare.lean
      -- thm:bv-poincare
    PoincareTrace.lean
      -- prop:poincare-trace
    RelativeCubes.lean
      -- prop:rel-iso-cube
    RelativeScaling.lean
      -- prop:rel-iso-ball
      -- prop:rel-iso-disk
    Rellich.lean
      -- thm:rellich

  Stationary/
    Bootstrap.lean
      -- lem:mc-elliptic
      -- prop:bootstrap-C2
      -- prop:bootstrap-C3
      -- cor:EL-pointwise
    CapacitaryEstimate.lean
      -- not:cap-estimate
      -- lem:integrate-EL
      -- lem:two-bounds-I
      -- lem:eliminate-capacity
      -- prop:cap-estimate
    Connected.lean
      -- prop:connected
    Defs.lean
      -- def:stationary
      -- not:stationary
    EulerLagrange.lean
      -- thm:constrained-first-var
      -- lem:graph-PMC
    MinimizerContext.lean
      -- not:minimizer-rep
    ScalingIdentity.lean
      -- prop:scaling-identity
      -- cor:minimizer-stationary

  Surface/
    Geometry.lean
      -- conv:curvature
      -- def:frame-free
    Morse.lean
      -- lem:morse-height
      -- lem:morse-distinct
      -- lem:morse-coords
      -- lem:one-manifold
      -- lem:morse-band
      -- cor:sublevel-stable
      -- lem:local-sectors
    MorseCount.lean
      -- def:merging-saddle
      -- lem:level-adjacency
      -- lem:sublevel-closure
      -- lem:count-submerging
      -- lem:merge-disjoint
      -- prop:morse-count
    Shape.lean
      -- lem:hess-height
      -- lem:gauss-jacobian
      -- conv:gauss-orientation
      -- lem:sign-kappa-index
      -- lem:critical-normal
    TotalCurvature.lean
      -- lem:regular-fiber-finite
      -- def:index-sum
      -- lem:gauss-map-extension
      -- thm:total-curvature-index
      -- prop:index-sum-antipodal
      -- thm:total-curvature-bound

  Tests/
    Constants.lean
      -- test:product-cost
      -- test:constants
    Curvature.lean
      -- test:total-curvature
      -- test:antipodal
    Iteration.lean
      -- test:recurrence
      -- test:excess-values
    LevelFunctionals.lean
      -- test:K-model
      -- test:K-bochner
    Signs.lean
      -- test:model-potential
      -- test:ball-stationary
      -- test:green-normal
      -- test:tubular-flux
      -- test:kappa-sign

  Threshold/
    Algebra.lean
      -- lem:Vstar-form
      -- lem:Vstar-identity
      -- lem:Vstar-bounds
    Comparison.lean
      -- prop:two-ball
    Defs.lean
      -- def:threshold
    Ledger.lean
      -- lem:ledger-Fprime
      -- lem:ledger-oneball
      -- lem:ledger-six
      -- lem:ledger-cubic
      -- lem:ledger-binding
      -- lem:ledger-L5
      -- lem:ledger-h

  Topology/
    Parity.lean
      -- prop:orientation-parity
      -- thm:component-count
    Transversality.lean
      -- thm:transversality
      -- cor:transv-preimage
      -- lem:handshake
    Tubular.lean
      -- thm:tubular

  Value/
    BoundedApprox.lean
      -- lem:bounded-approx
    Continuity.lean
      -- lem:m-concave
      -- prop:m-continuous
    Defs.lean
      -- def:value
      -- lem:m-finite
      -- lem:subadditivity

  Variation/
    ConstantTranslation.lean
      -- lem:sym-diff-constant
    CoulombBoundary.lean
      -- lem:coulomb-first-var
    FieldStability.lean
      -- lem:sym-diff-stability
    OneDimensionalTranslation.lean
      -- lem:1d-translation
    Perimeter.lean
      -- thm:first-var-perimeter
    Piola.lean
      -- lem:piola
    Straight.lean
      -- conv:first-variation
    StraightCofactor.lean
      -- lem:straight-cofactor
    StraightDiffeo.lean
      -- def:straight
      -- lem:straight-diffeo
    SymmetricDifferenceFull.lean
      -- thm:sym-diff
    TransportC1.lean
      -- thm:transport-perimeter
    VolumeBoundary.lean
      -- thm:first-var-volume
\end{verbatim}
}


\section{Notes on the tree}

\begin{enumerate}[leftmargin=*]

\item \textbf{\code{Elliptic/ClassicalCalculus.lean} must be built early.}
  Despite living under \texttt{Elliptic/}, it is imported by
  \texttt{DeGiorgi/Structure.lean} (for
  Corollary~\ref{cor:smooth-perimeter}).  It depends only on
  \texttt{BV/Rellich.lean} and chart calculus.  See
  Remark~\ref{rem:acyclic-elliptic}.  If a strict layering is preferred, move
  it to \code{Calculus/ChartGaussGreen.lean} at the top of the tree; the
  labels are unchanged.

\item \textbf{\code{Isoperimetric/Rigidity.lean} is deliberately far from
  \texttt{Isoperimetric/ABP.lean}.}  \texttt{ABP.lean} depends on
  \texttt{Elliptic/} only; \texttt{Rigidity.lean} depends on the whole
  \texttt{Regularity/} tree.  Keeping them in one file would create a false
  impression of a cycle and would force the whole regularity development to be
  imported by every consumer of \eqref{eq:sharp-isoperimetric}.  See
  Remark~\ref{rem:acyclic-isoperimetric}.

\item \textbf{\code{Regularity/OmegaMinimal.lean} holds only the
  definition.}  This lets \code{Isoperimetric/Rigidity.lean} and
  \code{Regularity/Penalization.lean} both supply instances without either
  importing the other, per (R\ref{r:generic}).

\item \textbf{The threshold algebra is the first pilot.}  Start with
  \texttt{Conventions.lean} and \texttt{Threshold/Defs.lean}, then
  \texttt{Threshold/Algebra.lean} and \texttt{Threshold/Ledger.lean}.
  The ledger is pure algebra, and its statements
  (Lemmas~\ref{lem:ledger-Fprime}--\ref{lem:ledger-h},
  \ref{lem:ledger-L5}) are exactly the ones a normalisation tactic can
  discharge.  Doing it first de-risks the sharp constants and gives an early
  win.  Note that $q=2^{1/3}$ should be introduced once, as a real cube root
  with $q^3=2$ and $q>1$, and never as a \texttt{Real.rpow} expression to be
  re-derived.  Prove agreement with the showcase's
  \texttt{LiquidDrop.criticalVolume}.  The two-ball comparison is assigned
  to \code{Threshold/Comparison.lean}, so the pilot does not import ball
  energy or its geometric prerequisites.

\item \textbf{\texttt{Flow/} plus \texttt{Surface/Morse.lean} is the
  optional-cost centre.}  \texttt{Flow/} exists only for
  \texttt{lem:morse-band} and \texttt{lem:one-manifold}, which exist only for
  \texttt{Surface/MorseCount.lean}.  Formalisation
  note~\ref{fnote:retire-flows} gives the two ways this block might shrink,
  neither of them verified.  \texttt{Surface/Shape.lean} and the proof of the
  index-integral formula are independent of \texttt{Flow/}.  The complete
  \code{Surface/TotalCurvature.lean} file, however, contains
  Theorem~\ref{thm:total-curvature-bound}, which imports
  Proposition~\ref{prop:morse-count} from \texttt{Surface/MorseCount.lean}
  and therefore depends transitively on \texttt{Flow/}.  Build the complete
  file after the Morse count; deferring Flow would require a further file
  split or an explicit conditional interface for that count.

\item \textbf{Four files were deleted relative to the first draft of this
  document}: \code{Topology/Triangulation.lean},
  \texttt{Surface/Frame.lean}, \texttt{Surface/Stokes.lean} and
  \texttt{Surface/GaussBonnet.lean}, replaced by
  \texttt{Surface/Shape.lean}, \texttt{Surface/Morse.lean},
  \texttt{Surface/MorseCount.lean} and
  \code{Surface/TotalCurvature.lean}.  The reason is Formalisation
  note~\ref{fnote:why-degree}: the deleted route needed abstract simplicial
  complexes with geometric realisation, which the pinned mathlib does not
  provide in any usable form (\code{AlgebraicTopology/Simplicial*} is
  category-theoretic simplicial sets;
  \code{Analysis/Convex/SimplicialComplex} is convex geometry).  Note that
  no \texttt{Topology/Degree.lean} is needed either: Formalisation
  note~\ref{fnote:not-a-degree} explains why the index sum is never shown to
  be a degree.

\item \textbf{\texttt{Regularity/} has seventeen planned files.}  The
  $\varepsilon$-regularity argument is where a monolithic file would be least
  maintainable: the smallness constants of
  Constants~\ref{const:height}, \ref{const:graph} and \ref{const:tilt} are
  chosen in three different files and must be threaded in the declared order.
  Making each choice a separate file's \emph{output} --- an existential
  statement, per (R\ref{r:constants}) --- is what keeps the order honest.

\item \textbf{\texttt{Tests/} should be built incrementally, not at the end.}
  Add each test once its actual prerequisites are available.  The completed
  map records these dependencies explicitly: for example,
  \texttt{test:constants} also checks the optimal ball volume and binding
  energy, so its complete implementation comes after those results.  Its
  elementary threshold checks may be extracted earlier during the pilot.

\item \textbf{Thematic overview of the build.}  Following
  Chapter~\ref{ch:dependencies}:
  \texttt{Conventions} $\to$ \texttt{Threshold/Ledger} $\to$
  \texttt{Measure} $\to$ \texttt{BV} $\to$ \texttt{Area} $\to$
  \code{Elliptic/ClassicalCalculus} $\to$ \texttt{Sobolev} $\to$
  \texttt{DeGiorgi} $\to$ \texttt{Variation} $\to$ \texttt{Elliptic} (rest)
  $\to$ \texttt{Sard} $\to$ \texttt{BV/SmoothApprox} $\to$
  \texttt{Isoperimetric/ABP} $\to$ \texttt{Energy}, \texttt{Ball},
  \texttt{Threshold} $\to$ \texttt{Value}, \texttt{Binding/Splitting} $\to$
  \texttt{Compactness}, \texttt{Existence} $\to$ \texttt{Regularity} $\to$
  \texttt{Cones} $\to$ \texttt{Isoperimetric/Rigidity} $\to$
  \texttt{Stationary} $\to$ \texttt{Flow}, \texttt{Topology},
  \texttt{Surface} $\to$ \texttt{Hull}, \texttt{Capacity} $\to$
  \texttt{CapacitaryK} $\to$ \texttt{Green} $\to$
  \code{Stationary/CapacitaryEstimate} $\to$ \texttt{Classification},
  \texttt{Nonexistence} $\to$ \texttt{Main} $\to$ \texttt{Binding/Main}.

  The \texttt{Flow}/\texttt{Topology}/\texttt{Surface} block is placed late
  because nothing before \texttt{Hull} needs the complete block; its
  elementary surface definitions are available earlier.  For actual
  scheduling, use a topological order of the module imports, which
  interleaves these thematic groups as required.

\item \textbf{Blueprint node names.}  The \LaTeX{} labels in this document are
  intended to be the blueprint node names verbatim.  When a Lean declaration
  is written for an item, add \verb|\lean{Decl.Name}| to that item and
  \verb|\leanok| when it compiles without \texttt{sorry}.  Do not rename
  labels: Chapter~\ref{ch:dependencies} and this chapter both reference them.

\item \textbf{What may never appear.}  No \texttt{axiom}, no
  \texttt{constant}, no \texttt{sorry} in a finished file, and no
  \texttt{theorem} whose statement cites a published paper.  If an item cannot
  be proved from its declared \verb|\uses| set plus
  Chapter~\ref{ch:base}, that is a defect in this document and should be
  reported rather than patched with a hypothesis.
\end{enumerate}

\section{Provenance}

The mathematical target is O.~Chodosh and M.~Gianocca, \emph{No compromise in
the liquid drop model}, arXiv:2608.11517.  The source paper uses
Bonacini--Cristoferi for regularity and Agostiniani--Mazzieri for capacitary
monotonicity; \emph{neither is assumed here} --- they are replaced by
Chapters~\ref{ch:penalization}--\ref{ch:no-singular} and
Chapter~\ref{ch:appK} respectively, and are used only as provenance and as
checks on constants.  The slicing idea of Chapter~\ref{ch:nonexistence}
originates with Frank--Killip--Nam; the compactness and existence context of
Chapters~\ref{ch:existence} and \ref{ch:binding} is related to work of
Frank--Lieb and Frank--Nam.  The arguments needed here are written out in
full.
